%lezione 25 del 21 maggio 2026 pag 59 appunti
%-----------------------------------------------------------------------------------------------------
%Ma poi che fui ai pié d'un colle giunto
%là dove terminava quella valle
%che m'avea di paura il cor compunto
%guardai in alto, vidi le sue spalle
%vesite già dei raggi del pianeta
%che mena dritto altrui per ogni calle.
\subsection{Disuguaglianze di Bell}
Si consideri il medesimo sistema entangled: lo spin lungo 3 direzioni $a,b,c$, supponendo che la misura
di ciascuna componente includa il risultato delle altre 2, può avere $2^3=8$ combinazioni, con 
probabilità $N_i$:
\begin{table}[H]
\centering
\begin{tabular}{|c|cccccccc|}
    \hline
    N. di ptc.&$N_1$&$N_2$&$N_3$&$N_4$&$N_5$&$N_6$&$N_7$&$N_8$\\ \hline
    $\vartheta_a$&+&+&+&+&-&-&-&-\\
    $\vartheta_b$&+&+&-&-&+&+&-&-\\
    $\vartheta_c$&+&-&+&-&+&-&+&-\\ \hline
\end{tabular}
\caption{\small Combinazioni dello spin lungo le tre direzioni per la particella $a$ con relative
probabilità.}
\end{table}
La tabella per $b$ avrà tutti i segni opposti.\\
Si consideri ora la probabilità di misurare lo spin della particella $a$ parallelo a $\vartheta_a$
e quello della particella $b$ parallelo a $\vartheta_b$: secondo l'intepretazione di EPR vale
$$\mathcal{P}(\vartheta_a+,\vartheta_b+)=\frac{N_3+N_4}{2}$$
Similmente valgono
\begin{gather*}
\mathcal{P}(\vartheta_a+,\vartheta_c+)=\frac{N_2+N_4}{2}\\
\mathcal{P}(\vartheta_c+,\vartheta_b+)=\frac{N_3+N_7}{2}\\
\end{gather*}
Poiché $N_i\in[0,1]$, vale che
$$N_3+N_4\leq N_3+N_4+N_2+N_1 \leftrightarrow
\boxed{\mathcal{P}(\vartheta_a+,\vartheta_b+)\leq\mathcal{P}(\vartheta_a+,\vartheta_c+)+
\mathcal{P}(\vartheta_c+,\vartheta_b+)}$$
Questa è una formulazione della \textbf{disuguaglianza di Bell} ed essa è verificata per teorie
\textit{locali a variabili nascoste}: se le previsioni della Meccanica Quantistica soddisfano questa
disuguaglianza, allora EPR hanno ragione ed esistono gli elementi di realtà.

Le previsioni della Meccanica Quantistica sono quelle già ricordate nel paragrafo precedente:
la disuguaglianza di Bell è dunque
$$\sin^2\left(\frac{\vartheta_a-\vartheta_b}{2}\right)\leq\sin^2\left(\frac{\vartheta_a-\vartheta_c}{2}
\right)+\sin^2\left(\frac{\vartheta_c-\vartheta_b}{2}\right)$$
se si pone $\vartheta_a-\vartheta_b=4\alpha$ e $\vartheta_c-\vartheta_b=2\alpha$, la disuguaglianza
diventa $\cos(\alpha)\leq\sin(\alpha)$, che è \textbf{falsa} $\forall\alpha<\frac{\pi}{4}$:
questo \textbf{garantisce che la Meccanica Quantistica \textit{non} è una teoria locale a variabili
nascoste}.\\

Un'altra formulazione della disuguaglianza di Bell è la seguente: si supponga che in un sistema
ci siano $n$ variabili nascoste, rappresentate dal vettore $\mathbf{\lambda}\in\mathbb{R}^n$ e siano
$A(a,\mathbf{\lambda})$, $B(b,\mathbf{\lambda})$ i risultati di due misure, funzioni delle
caratteristiche del sistema e del rivelatore ($a$, $b$) e delle variabili nascoste
($\mathbf{\lambda}$); cambiando le caratteristiche dei rivelatori (l'angolo di misura per lo spin,
per esempio) i risultati diventano $A'(a',\mathbf{\lambda})$, $B'(b',\mathbf{\lambda})$.

Dopo aver effettuato le misure, si consideri la quantità
$$M(\mathbf{\lambda}):=AB-AB'+A'B+A'B'=A(B-B')+A'(B+B')$$
Se i possibili risultati delle misure sono solo $A,A',B,B'=\pm1$ si ha che
$$-2\leq \langle M\rangle\leq2$$
Nel caso dello spin, $A$ rappresenta la misura sulla particella $a$ in direzione $\vartheta_a$ e
similmente per $b$: si ha dunque che
\begin{gather*}
\langle AB\rangle=\langle\mathcal{P}_{\vartheta_a}\mathcal{P}_{\vartheta_a}\rangle=
(1)(1)\mathcal{P}_{++}+(1)(-1)\mathcal{P}_{+-}+(-1)(1)\mathcal{P}_{-+}+(-1)(-1)\mathcal{P}_{--}=\\
=\sin^2\left(\frac{\vartheta_a-\vartheta_b}{2}\right)-\sin^2\left(\frac{\vartheta_a-\vartheta_b}{2}
\right)=-\cos(\vartheta_a-\vartheta_b)
\end{gather*}

